% Source-derived chapter 02: Bialynicki-Birula partition and decomposition
% Original source: very-stable-higgs-bundles-equivariant-multiplicity-mirror-symmetry
\chapter{Bialynicki-Birula partition and decomposition}
\label{very-stable-higgs-bundles-equivariant-multiplicity-mirror-symmetry-chapter-02}
\source{very-stable-higgs-bundles-equivariant-multiplicity-mirror-symmetry}

\begin{remark}[Source section]
\label{very-stable-higgs-bundles-equivariant-multiplicity-mirror-symmetry-chapter-02-source}
\source{very-stable-higgs-bundles-equivariant-multiplicity-mirror-symmetry}
\source{very-stable-higgs-bundles-equivariant-multiplicity-mirror-symmetry}
This chapter reproduces the corresponding section of the retrieved
LaTeX source, including its definitions, statements, examples, and
proofs. It has not yet been translated into Lean.
\notready
\end{remark}

% The source section title was: Bialynicki-Birula partition and decomposition
\label{bb}
\source{very-stable-higgs-bundles-equivariant-multiplicity-mirror-symmetry}

\subsection{Bialynicki-Birula theory for semi-projective varieties}

Let $M$ be a normal, complex, semi-projective algebraic variety.  Semi-projective means that $M$ is quasi-projective and we have a $\T:=\C^\times$ action on $M$  such that  the fixed point locus $M^{\T}$ is projective and   for every $x\in M$ there is a $p\in M^\T$ such that
$\lim_{\lambda\to 0}  \lambda x=p$. By the latter we mean that there exists a $\T$-equivariant morphism $g:\A^1\to M$ with $g(1)=x$ and $g(0)=p$, where $\T$ acts on $\A^1$ in the standard way. 

Examples of semi-projective varieties include  cotangent bundles of smooth projective varieties, moduli spaces of Higgs bundles and Nakajima quiver varieties.

For every $\alpha\in M^{\T}$ denote by  $$W^+_\alpha:= \{ x\in M | \lim_{\lambda\to 0}  \lambda x = \alpha\}$$ the {\em upward flow} from $\alpha$ and $$W^-_\alpha:= \{ x\in M | \lim_{\lambda\to \infty}  \lambda x = \alpha\}$$ the {\em downward flow} from $\alpha$. We also define for connected components $F\in \pi_0(M^{\T})$ of the fixed point locus $$W^+_F:=\cup_{\alpha\in F} W^+_\alpha$$ the {\em attractor} of $F$  
and $$W^-_F:=\cup_{\alpha\in F} W^-_\alpha.$$ the {\em repeller} of $F$. 

Let $M^s\subset M$ denote the smooth locus of $M$. Let $\alpha\in M^{s\T}$ be a smooth fixed point.  Then $\T$ acts linearly on the tangent space $T_\alpha M$. We can decompose this representation into weight spaces
$T_\alpha M\cong \oplus_{k\in \Z} (T_\alpha M)_k$, where $(T_\alpha M)_k<T_\alpha M$ denotes the isotypical component, where $\lambda\in \T$ acts as multiplication by $\lambda^k$. We will denote by $T^+_\alpha M = \oplus_{k>0}(T_\alpha M)_k$ the positive and by $T^-_\alpha M = \oplus_{k<0}(T_\alpha M)_k$ the negative part of this decomposition. Let  $T^0_\alpha M=(T_\alpha M)_0$ denote the zero weight component then we have the decomposition $$T_\alpha M \cong  T^+_\alpha M \oplus T^0_\alpha M \oplus T^-_\alpha M. $$

\begin{proposition}
\source{very-stable-higgs-bundles-equivariant-multiplicity-mirror-symmetry}
\notready\label{bbprop} For a smooth fixed point $\alpha\in M^{s\T}$ the upward and downward flows  $W^+_\alpha$ and $W^-_\alpha$ are $\T$-invariant locally closed subvarieties in $M$. Furthermore we have $W^+_\alpha \cong T^+_\alpha M$ and $W^-_\alpha\cong T^-_\alpha M$ as varieties with $\T$-action.  \end{proposition}
\source{very-stable-higgs-bundles-equivariant-multiplicity-mirror-symmetry}
\begin{proof} This is proved in \cite{bialynicki-birula} for a smooth complete $M$. The statement above could be reduced to that case by equivariantly compactifying $M$ (see below or \cite[Theorem 3]{sumihiro}) and equivariantly resolving the singularities of the compactification (\cite[Corollary 7.6.3]{villamayor}). 
	
	For completeness, and to prepare the ground for the proof of Proposition~\ref{lagrangian}, we present a proof  inspired by the approach of \cite[Lemma 2.2]{bellamy-etal} see also \cite[\S 1]{drinfeld3}.  
	Let $\alpha\in M^{s\T}$.  As the singular locus $\operatorname{Sing}(M)\subset M$ is $\T$-invariant and closed we have   $W^+_\alpha\subset M^s$. Take $\alpha\in U_0\subset M^s$  a $\T$-invariant open affine neighbourhood (using \cite[Corollary 2]{sumihiro}). In particular, $U_0$ is smooth.  The $\T$-action on $U_0$ induces a $\Z$-grading on $\C[U_0]\cong H^0(U_0,\calO_{U_0})$, defined by setting $f\in \C[U_0]_i$ if and only if \beq\label{transrule}f(\lambda p)=\lambda^if(p)\eeq for all $\lambda \in \T$ and $p\in U_0$. In this case we will say that $f$ is homogeneous of degree $\deg(f)=i$. 
\source{very-stable-higgs-bundles-equivariant-multiplicity-mirror-symmetry}
	
	Let $\m_\alpha\triangleleft \C[U_0]$ be the maximal ideal of functions vanishing at $\alpha$. As $\alpha$ is $\T$-fixed $\m_\alpha$ is a homogeneous ideal. We can find homogeneous elements $x_1,\dots,x_{n}\in \m_\alpha\subset \C[U_0]$, where $n=\dim(M)$ such that their image  in $\m_\alpha/\m_\alpha^2\cong T^*_\alpha M$ form a basis. The one-forms $dx_1,\dots, dx_n\subset H^0(U_0,\Omega_{U_0})$ then form a basis at $\alpha$ and we let $U$ be the complement of the divisor of zeroes of the section $dx_1\wedge\dots\wedge dx_n\in H^0(U_0,\Omega^n(U_0))$. It is a $\T$-invariant affine open subset $U\subset U_0$, where $dx_1,\dots, dx_n$ remains a basis at every point. In particular, $U$ is non-singular. 
	As $U$ is invariant under the $\T$-action and the closure of any $\T$-orbit in $\M\setminus U$ is still contained in $\M\setminus U$, as it is closed,  we have   $$W^{\pm}_\alpha \subset U.$$ Therefore we can prove Proposition~\ref{bbprop} only working in $U$.
	%We can also assume that  $\{\alpha\}=V(x_1,\dots,x_d)\subset U$ that is all the $x_i$ only vanish at $\alpha$ . % Let $\{1,\dots,n\}=U\coprod \calF\coprod {\mathcal D}, $ such that $\deg(x_i)>0$ when $i\inU$, $\deg(x_i)=0$, when $i\in \calF$ and $\deg(x_i)<0$ when $i\in {\mathcal D}$.
	
	Let $A:=\C[U]$ and $A=\oplus_{n\in \Z} A_n$ be the grading induced by the $\T$-action. Let $A_+=\oplus_{n\in \Z_{>0}} A_n$ and $A_-=\oplus_{n\in \Z_{<0}} A_n$. Thus in particular $A=A_-\oplus A_0\oplus A_+$.  First we observe that 
	
	\begin{lemma}
\source{very-stable-higgs-bundles-equivariant-multiplicity-mirror-symmetry}
\notready $U^\T=V(f|f\in A_++A_-)\subset U$ i.e. the subvariety of zeroes of homogenous $f$ of non-zero degree. Thus $\C[U^\T]=A/(A_+,A_-)$. \end{lemma}
	\begin{proof} Clearly if $p\in U^\T$, $f \in A_i$ with $i\in \Z\setminus \{0\}$    and $\lambda\in \T$ such that $\lambda^i\neq 1$ then $$f(p)=f(\lambda p)=\lambda^if(p)=0.$$ On the other hand if $p\in U$ is such that  $f(p) =0$  for all $f\in A_i$ with $i\in \Z\setminus \{0\}$ then the maximal ideal at $p$ satisfies $(A_++A_-)\subset \m_p$  and so $\m_p$ is a homogenous ideal, thus $p\in U^\T$. \end{proof}
	
	Let $F:=U^\T$ and, as above, $W^+_{F}=\coprod_{\beta\in F} W^+_\beta=\{p\in U|\lim_{\lambda\to 0}\lambda p\in F\}$ the attractor of $F$ in $U$. Then we have 
	\begin{lemma}
\source{very-stable-higgs-bundles-equivariant-multiplicity-mirror-symmetry}
\notready \label{ringofw} $W^+_{F}=V(f|f\in A_-)$ i.e. the subvariety of zeroes of homogenous $f$ of negative degree. Thus $\C[W^+_{F}]=A/(A_-)$. 
\source{very-stable-higgs-bundles-equivariant-multiplicity-mirror-symmetry}
	\end{lemma}
	\begin{proof} Again it is clear that if $p\in W^+_{F}$ then $f(p)=0$ for $f\in A_-$. Namely if $f\in A_i$ for $i<0$ then $$f(\lim_{\lambda\to 0} \lambda p)=\lim_{\lambda\to 0} f(\lambda p)=\lim_{\lambda\to 0} \lambda^{i} f(p)$$ implies that $f(p)=0$.
		
		On the other hand assume that for some $p\in U$ we have  $f(p)=0$ for all $f\in A_-$. Then let $\pi_p:A\to A/\m_p\cong \C$ and define \bes \begin{array}{cccc}g:&A&\to& \C[x]\\ &\sum_i a_i&\mapsto& \sum_i \pi_p(a_i)x^i
		\end{array},\ees where we write every element in $A$ in the form of a finite sum of homogeneous elements $\sum_i a_i$ where $a_i\in A_i$. Because $f(a_i)=0$ for $i<0$, the definition makes sense. We see that $g$ is a graded ring homomorphism.  Observe also that for $z\in \A^1=\Spec(\C[x])$ the composition $$(\pi_{z}\circ g)\left(\sum_i a_i\right)=\sum_i \pi_p(a_i)z^i.$$ When $z= 1$ we see that $\ker(\pi_{1}\circ g)=\m_{p}$ and that $\ker(\pi_0\circ g)=(\ker(\pi_p|_{A_0}),A_-,A_+)$ is a maximal ideal containing $A_-$ and $A_+$ thus corresponds to a fixed point of the $\T$-action. Indeed the ring homomorphism $g$ gives a $\T$-equivariant map $\A^1\to U$, such that $g(1)=p$ and $g(0)=\lim_{\lambda\to 0} \lambda p\in F$ by  definition. Thus $p\in W^+_{F}$. The result follows.
	\end{proof}
	
	\begin{remark}
\source{very-stable-higgs-bundles-equivariant-multiplicity-mirror-symmetry}
\notready We can build a map $\pi:W^+_{F}\to F$ by sending $p\in W^+_{F}$ to $\lim_{\lambda \to 0} \lambda p\in F$. By the construction above we see that the maximal ideal of $\pi(p)$ is $(\ker(\pi_p|_{A_0}),A_-,A_+)\triangleleft A$. Thus we see that $\pi$ is induced by the map \beq\label{pistar}\begin{array}{cccc}\pi^*:&A/(A_+,A_-)&\to& A/(A_-)\\ &a \mbox{ \rm mod } (A_+,A_-)&\mapsto & a_0 \mbox{ mod } (A_-)\end{array},\eeq where $a=\sum_{i\in \Z} a_i$ is the sum of homogeneous elements $a_i\in A_i$. It is straightforward to check that $\pi^*$ is a well-defined graded ring homomorphism and that $(\pi^*)^{-1}(\ker(\pi_p))=(\ker(\pi_p|_{A_0}),A_-,A_+)$.   
\source{very-stable-higgs-bundles-equivariant-multiplicity-mirror-symmetry}
	\end{remark}
	
	Next we determine the tangent spaces of $F$, $W^+_{F}$, $W^-{F}$, $W^+_\alpha$ and $W^-_\alpha$. To do this we will denote by $\partial_i\in H^0(U,TU)$ the dual basis to $dx_i$. Recall that  $\lambda\in \T$ acts on the function $x_i\in\C[U]$ via the formula $$\lambda\cdot x_i(p)=x_i( \lambda^{-1}p)=\lambda^{-\deg(x_i)}x_i(p),$$ i.e. $\lambda\cdot x_i=\lambda^{-deg(x_i)}x_i$ which should not be confused with the transformation rule in \eqref{transrule}.  In particular, the induced action on one-forms gives \beq\label{transdxi} \lambda\cdot d x_i =(m_\lambda)_*(d x_i)=\lambda^{-deg(x_i)}dx_i, \eeq where $m_\lambda$ denotes the action isomorphisms:  $$\begin{array}{cccc}m_\lambda: &U&\to& U\\ &p&\mapsto& \lambda p\end{array}$$ for $\lambda\in \T$. One can compare\footnote{There seems to be a related confusion in \cite[Definition 1.1]{bellamy-etal}. Their  definition of elliptic $\T$-action  should either require the existence of limit points in the $t\to 0$ limit as in this paper, or replace the definition of positive weight using $(m_t)_*$ instead of $m_t^*$. In fact, when they use the definition of positive weight in the displayed line after (2.1) on page 2613 of \cite{bellamy-etal}, they use the correct $(m_t)_*$ instead of $m_t^*$ as in   \cite[Definition 1.1]{bellamy-etal}. } \eqref{transdxi} to $$m_\lambda^*(dx_i)=d(m_\lambda^*(x_i))=\lambda^{\deg(x_i)}dx_i .$$  This way we see that $\partial_i$ will have weight $\deg(x_i)$ in the sense that \beq\label{homdi} (m_\lambda)_*(\partial_i)=T m_\lambda(\partial_i)=\lambda^{\deg(x_i)}(\partial_i).\eeq  As a reality check for $\lambda\in \T$  we compute, using \eqref{transrule} $$dx_j(T m_\lambda (\partial_i))=m_\lambda^*(dx_j)(\partial_i)=d ( x_j\circ m_\lambda)(\partial_i)=  \lambda^{\deg x_j} dx_j(\partial_i)=\lambda^{\deg x_j} \delta_{i,j},$$ showing \eqref{homdi}.
\source{very-stable-higgs-bundles-equivariant-multiplicity-mirror-symmetry}
	\begin{lemma}
\source{very-stable-higgs-bundles-equivariant-multiplicity-mirror-symmetry}
\notready For $\beta\in F$ we have $$T_\beta F = (T_\beta U)^\T\cong \operatorname{span}(\partial_i|_\beta|\deg(x_i)=0)\subset T_\beta U.$$ In particular, $F\subset U$ is a non-singular subvariety. 
	\end{lemma}
	\begin{proof} Dually, we will prove that the kernel of the surjective map $g:T^*_\beta U\to T^*_\beta F$ is $$\operatorname{span}(dx_i|_\beta| \deg(x_i)\neq 0)<T^*_\beta U.$$ It is clear that $dx_i|_\beta\in \ker(g)$ when $\deg(x_i)\neq 0$, because $g$ is a map of $\T$-modules when we endow $T^*_\beta F$ with the trivial $\T$-module structure. Let $\eta\in (T^*_\beta U)^\T\cap \ker(g)$. We can represent $\eta=d {y}_0$ with the homogeneous degree $0$ element ${y}_0\in \mathfrak{m}_{\beta,U}\triangleleft A$. We will  denote by $$\bar{y}_0\in \mathfrak{m}_{\beta,F}\triangleleft A/(A+,A_-)\cong \C[F]$$ the image of ${y}_0$ in the quotient $A/(A+,A_-)$.  By assumption $0=g(\eta)=d{\bar{y}}_0$ thus $\bar{y}_0\in \mathfrak{m}_{\beta,F}^2.$ Thus there exists some $\bar{c}_i,\bar{d}_i\in  \mathfrak{m}_{\beta,F}$ such that $\bar{y}_0=\sum_i \bar{c}_i\bar{d}_i$. Denoting by ${c}_i$ and ${d}_i$ some lift of $\bar{c}_i$ and $\bar{d}_i$ in $\mathfrak{m}_{\beta,U}$, we get that ${y_0}-\sum_i {c}_i{d}_i\in (A_+,A_-)$. We can assume that  $c_i$ and $d_i$ are homogeneous and in turn that $\deg(c_i)=\deg(d_i)=0$.   We have  some $a_j\in A_+$, $b_j\in A_-$ such that ${y_0}-\sum_i {c}_i{d}_i=\sum_j a_jb_j$. We can assume that $a_j$ and $b_j$ are homogeneous of non-zero degree and in turn that $\deg(a_j)+\deg(b_j)=0$. 
		Denoting $\pi_\beta(a_j)=\lambda_j \in A/\m_{\beta,U}\cong \C$ and $\pi_\beta(b_j)=\mu_j \in A/\m_{\beta,U}\cong \C$ we get that $a_j-\lambda_j$ and $b_j-\mu_j$ are in $\m_{\beta,U}$. Thus we  conclude that $$y_0-(\sum_j a_j \mu_j+\lambda_j b_j)\in \m_{\beta,U}^2.$$ As $\m_{\beta,U}^2$ is a homogeneous ideal we deduce that $y_0\in \m_{\beta,U}^2$ proving the first claim. 
		
		Finally, we deduce that $\dim(\m_{\beta,U}/\m^2_{\beta,U})=\#\{dx_i|\deg(x_i)=0\}$ is independent of $\beta\in F$, thus indeed $F$ is non-singular.
		% $\mathfrak{m}^2_{\beta,F}\triangleleft \C[F]\cong A/(A_+,A_-)$ is the maximal ideal corresponding to the point $\beta\in F$. 
	\end{proof}
	\begin{lemma}
\source{very-stable-higgs-bundles-equivariant-multiplicity-mirror-symmetry}
\notready \label{tangentattractor} For $p\in W^+_{F}$ we have $T_p W^+_{F}=\operatorname{span}(\partial_i|_p|\deg(x_i)\geq 0)<T_pU.$ Similarly for $p\in W^-_{F}$ we have $T_p W^-_{F}=\operatorname{span}(\partial_i|_p|\deg(x_i)\leq 0)<T_pU.$
\source{very-stable-higgs-bundles-equivariant-multiplicity-mirror-symmetry}
	\end{lemma}
	\begin{proof} Again, we will prove dually that the kernel of the surjective map $h:T^*_p U\to T^*_p W^+_{F}$ is $$\operatorname{span}(dx_i|_p| \deg(x_i)< 0)<T^*_p U.$$ By Lemma~\ref{ringofw} for $\deg(x_i)<0$ the function $x_i$ vanishes on $W^+_{F}$ and thus $dx_i|_p$ is in the kernel of $h$. 
		Assume now that $\sum_{\deg x_i\geq 0} \alpha_i dx_i|_p\in \ker(h)$ for some $\alpha_i\in \C$.  This means that $$\sum_{\deg x_i\geq 0} \alpha_i (x_i-\pi_p(x_i))\in \m_{p,W^+_{F}}^2.$$ It follows that there are $c_j,d_j\in \m_{p,U}$ such that $$\sum_{\deg x_i\geq 0} \alpha_i (x_i-\pi_p(x_i))-\sum_j c_j d_j\in (A_-).$$ So we have homogeneous $a_k\in A_-$ and $r_k\in A$ such that $$\sum_{\deg x_i\geq 0} \alpha_i (x_i-\pi_p(x_i))-\sum_k r_k a_k\in \m_{p,W^+_{F}}^2 .$$ By writing $r_k=r_k-\pi_p(r_k)+\pi_p(r_k)$ and noting that $\pi_p(a_k)=0$ as $p\in W^+_{F}$ and $a_k\in A_-$ we get
		$$\sum_{\deg x_i\geq 0} \alpha_i (x_i-\pi_p(x_i))-\sum_k \beta_k a_k\in \m_{p,W^+_{F}}^2 ,$$ where $\beta_k=\pi_p(r_k)\in \C$. In other words \beq \label{form0} 0=\sum_{\deg x_i\geq 0} \alpha_i dx_i|_p-\sum_k \beta_k da_k|_p \in T^*_p(U).\eeq Our final observation is that for $\deg(x_i)\geq 0$ the function $\partial_i(a_k)$ has degree $-\deg(x_i)+\deg(a_k)<0$. This is because $da_k=  \sum_j \partial_j(a_k)dx_j$ therefore $$ da_k\wedge dx_1\wedge\dots \wedge dx_{i-1}\wedge d x_{i+1}\wedge \dots \wedge dx_n=(-1)^{i-1} \partial_i(a_k)dx_1\wedge \dots \wedge dx_n. $$ Thus $\partial_i(a_k)$ vanishes on $W^+_{F}$. This means that $da_k|_p$ is a linear combination of the $dx_i$ with $\deg(x_i)<0$ and so $\sum_{\deg x_i\geq 0} \alpha_i dx_i|_p$ and $\sum_k \beta_k da_k|_p$ are linearly independent. Thus \eqref{form0} implies that $\sum_{\deg x_i\geq 0} \alpha_i dx_i|_p=0$. The first statement follows. 
\source{very-stable-higgs-bundles-equivariant-multiplicity-mirror-symmetry}
		
		The proof of the second statement on the downward flows is  the same as above for the inverse action of $\T$ on $U$.   \end{proof}
	
	\begin{lemma}
\source{very-stable-higgs-bundles-equivariant-multiplicity-mirror-symmetry}
\notready \label{tangentup} For $\alpha\in F$, and $p\in W^+_\alpha$  we have  $T_p W^+_{\alpha}=\operatorname{span}(\partial_i|_p|\deg(x_i)> 0)<T_pU.$ Similarly,  for $p\in W^-_\alpha$  we have  $T_p W^-_{\alpha}=\operatorname{span}(\partial_i|_p|\deg(x_i)< 0)<T_pU.$
\source{very-stable-higgs-bundles-equivariant-multiplicity-mirror-symmetry}
	\end{lemma}
	
	\begin{proof} First we note that the locus of zeroes of $(x_1,\dots,x_n)$ in $U$ is zero dimensional, because $dx_1,\dots,dx_n$ is a basis at every point. Also this locus is invariant by $\T$, they are thus fixed points in $F$. All non-negative degree elements vanish on $F$ therefore the zeroes of the $0$ degree elements $${\bf x}_0:=\{x_i|\deg(x_i)=0\}$$ on $F$ are precisely the zeroes of $(x_1,\dots,x_n)$ on $U$. 
		
		By construction one such common zero is $\alpha$ and so $W^+_\alpha$ is a connected component of the variety of zeroes of $\bf{x}_0$ on $W^+_{F}$. In particular, $W^+_\alpha\subset U$ is closed and thus it is locally closed in $M$. 
		
		Thus we have that  $W^+_\alpha$ is a connected component of $\Spec(A/(A_-,\bf{x}_0))$. For $p\in W^+_\alpha$ we determine the kernel of the surjective map $f:T^*_p U\to T^*_p W^+_\alpha$. All functions in $A_-$ and  $\bf{x}_0$ restrict to $W^+_\alpha$ as zero thus $dx_i|_p$ is in the kernel when $\deg(x_i)\leq 0$. Let now $\sum_{\deg x_i>0} \alpha_i dx_i|_p \in \ker(f)$. This means that $$\sum_{\deg x_i> 0} \alpha_i (x_i-\pi_p(x_i))\in \m_{p,W^+_\alpha }^2,$$ hence there are $c_j,d_j\in \m_{p,U}$ such that $$\sum_{\deg x_i> 0} \alpha_i (x_i-\pi_p(x_i))-\sum_j c_j d_j\in (A_-,\bf{x}_0).$$ In turn, this means that there are homogeneous elements $a_k\in A_-$  and scalars $\beta_j\in \C$ such that \beq \label{final0} 0=\sum_{\deg x_i>0} \alpha_i dx_i|_p-\sum_k  da_k|_p-\sum_{\deg x_j=0} \beta_j dx_j|_p\in  T^*_p(U).\eeq Just as in the proof of Lemma~\ref{tangentattractor} above we can argue that as $a(p)=0$ for all $a\in A_-$, $da_k|_p\in \operatorname{span}(dx_i|_p|\deg x_i<0)$, thus the linear independence of $\{dx_1|_p,\dots,dx_n|_p\}$ and \eqref{final0} implies that $\sum_{\deg x_i>0} \alpha_i dx_i|_p =0$ proving the first claim.
\source{very-stable-higgs-bundles-equivariant-multiplicity-mirror-symmetry}
		
		The second claim again follows from the first by inverting the $\T$-action on $U$.
	\end{proof}
	
	To finish the proof of Proposition~\ref{bbprop} we define the map \beq\label{model}f_\alpha:W^+_\alpha\to T^+_\alpha\eeq sending $$p\mapsto \sum_{\deg(x_i)>0} x_i(p)\partial_i$$ where $\partial_i\in T_\alpha M$ is the dual basis of $dx_1,\dots,dx_n$ at $\alpha$. By Lemma~\ref{tangentup} the derivative of $f_\alpha$  is an isomorphism, thus it is \'etale and so open. Furthermore $f_\alpha$ is $\T$-equivariant and  its image is an open $\T$-invariant subvariety in the positive $\T$-module $T^+_\alpha M$, containing the origin, thus $f_\alpha$ is surjective. As $T^+_\alpha M$ is simply connected, it follows that $f_\alpha$ is a trivial covering, and $W^+_\alpha$ being connected implies that $f_\alpha$ is an isomorphism. 
\source{very-stable-higgs-bundles-equivariant-multiplicity-mirror-symmetry}
\end{proof}

\begin{remark}
\source{very-stable-higgs-bundles-equivariant-multiplicity-mirror-symmetry}
\notready \label{afffibr}
\source{very-stable-higgs-bundles-equivariant-multiplicity-mirror-symmetry}
	If we put the maps $f_\alpha$ in a family over $F$ we get a $\T$-equivariant isomorphism $f:W^+_{F}\cong  T^+F$ showing that $U^s_{F_\alpha}\subset M$ is a locally trivial affine fibration, another main result of \cite{bialynicki-birula}.  
\end{remark}


\begin{definition}
\source{very-stable-higgs-bundles-equivariant-multiplicity-mirror-symmetry}
\notready	We call $M=\coprod _{\alpha}W^+_\alpha$ the {\em Bialynicki-Birula partition} and $M=\coprod _{F}W^+_F$ the {\em Bialynicki-Birula decomposition}. Define also \beq\label{core}\calC:=\coprod_{\alpha\in M^\T} W^-_\alpha=\coprod_{F\in \pi_0(M^\T)} W^-_F\eeq to be the {\em core} of $M$. 
\source{very-stable-higgs-bundles-equivariant-multiplicity-mirror-symmetry}
	
\end{definition}

Additionally, in all our examples  $(M^s,\omega)$ with $\omega\in \Omega^2(M)$ will be  symplectic  s.t. \beq \label{weightone}\lambda^*(\omega)=\lambda\omega.\eeq Then we have
\source{very-stable-higgs-bundles-equivariant-multiplicity-mirror-symmetry}
\begin{proposition}
\source{very-stable-higgs-bundles-equivariant-multiplicity-mirror-symmetry}
\notready \label{lagrangian} Let $M$ be a normal, semi-projective complex variety with $(M^s,\omega)$ symplectic with $\omega\in \Omega^2(M^s)$. Assume \eqref{weightone} that $\omega$ is homogeneous of weight $1$.  For a smooth point $\alpha\in M^{s\T}$ the subspaces $T^+_\alpha M, T^{\leq 0}_\alpha M:=T^0_\alpha M\oplus T^+_\alpha M <T_\alpha M$ and subvarieties $W^+_\alpha,W^-_{F_\alpha}\subset M^s$ are Lagrangian.
\source{very-stable-higgs-bundles-equivariant-multiplicity-mirror-symmetry}
\end{proposition}
\begin{proof} As $\alpha\in M^{s\T}$ the tangent space $T_\alpha M$ is a $\T$-module. Let $X,Y\in T_\alpha X$ be homogeneous  with weight $\nu_1$ and $\nu_2$ respectively. That means that $\lambda\cdot X=\lambda^{\nu_1} X$ and $\lambda\cdot X=\lambda^{\nu_2} Y$. Then from \eqref{weightone} $$\lambda\omega(X,Y)=\lambda^*(\omega)(X,Y)=\omega(\lambda\cdot X,\lambda\cdot Y)=\lambda^{\nu_1+\nu_2}\omega(X,Y).$$ Consequently $\omega(X,Y)=0$ unless $\nu_1+\nu_2=1$. Thus $\omega$ is trivial on $T^+_\alpha$ and $T^{\leq 0}_\alpha$. As $T_\alpha M=T^+_\alpha\oplus T^{\leq 0}_\alpha$ both subspaces are Lagrangian.
	
	Recall the construction of $U\subset M^s$, the $\T$-invariant affine neighbourhood of $\alpha\in U$ from the proof of Proposition~\ref{bbprop} above.  Let $\omega=\sum_{i<j} f_{i,j} dx_i\wedge dx_j$ for some unique $f_{i,j}\in \C[U]$. As $\omega$, $x_i$ and $x_j$ are homogeneous so is $f_{i,j}$ and $\deg{f_{i,j}}=1-\deg(x_i)-\deg(x_j)$. Let now $p\in W^+_\alpha $. By Lemma~\ref{tangentup} the tangent vectors $\partial_k|_p\in T_p U$ with positive weights will be a basis for $T_p W^+_\alpha$. Take two such $\partial_k|_p, \partial_l|_p\in T_pW^+_\alpha$ and compute $$\omega(\partial_k|_p,\partial_l|_p)=\sum_{i<j} f_{i,j} dx_i\wedge dx_j(\partial_k|_p,\partial_l|_p)=f_{k,l}(p).$$
	The degree of $x_l$ and $x_k$ are positive, so the degree of $f_{k,l}$ is negative. In particular, $f_{k,l}$ vanishes on $W^+_\alpha$. Thus $f_{k,l}(p)=0$, and $T_pW^+_\alpha\subset T_p U$ is isotropic, and half-dimensional as $T_\alpha W^+_\alpha=T_\alpha^+<T_\alpha M$ is Lagrangian. Thus $W^+_\alpha$ is indeed Lagrangian in $U\subset M^s$. Similarly we can prove that $W^-_F$ is Lagrangian in $M^s$. The result follows.\end{proof}
\begin{remark}
\source{very-stable-higgs-bundles-equivariant-multiplicity-mirror-symmetry}
\notready Our approach in the  proof of Proposition~\ref{lagrangian} above was inspired by \cite[\S 2.2]{bellamy-etal}. Similar results were mentioned in \cite[Proposition 7.1]{nakajima}.
\end{remark}
\subsection{Very stable upward flows}
\begin{definition}
\source{very-stable-higgs-bundles-equivariant-multiplicity-mirror-symmetry}
\notready Let $\alpha\in M^{s\T}$. We say that $\alpha$ is {\em very stable} if the only point in its upward flow which is also in the core $\calC$ of \eqref{core} is $\alpha$: $$W^+_\alpha\cap \calC=\{\alpha\}.$$ % We also say that $\alpha\in M^{s\T}$ is {\em wobbly} if it is not very stable. 
\end{definition} 

\begin{remark}
\source{very-stable-higgs-bundles-equivariant-multiplicity-mirror-symmetry}
\notready The origin of the term {\em very stable} is in the work of Drinfeld \cite{drinfeld} and  Laumon \cite{laumon}, which was applied to vector bundles on curves (see Section~\ref{verystable}). However there is an accidental coincidence in  terminology. Namely, Bialynicki-Birula \cite{bialynicki-birula2} using terminology of Smale \cite[II.2]{smale} calls our upward flow $W^+_\alpha$ the {\em stable} subscheme of $\alpha$ while our downward flow $W^-_\alpha$ is called the {\em unstable} subscheme of $\alpha$. Our definition above can be formulated to say that the stable subscheme $W^+_\alpha$ is {\em very stable} if and only if it disjoint from the unstable subschemes $W^-_\beta$ for $\beta\neq \alpha$. 
\end{remark}
\begin{proposition}
\source{very-stable-higgs-bundles-equivariant-multiplicity-mirror-symmetry}
\notready \label{upwardclose} $\alpha\in M^{s\T}$ is very stable if and only if $W^+_\alpha\subset M$ is closed. 
\source{very-stable-higgs-bundles-equivariant-multiplicity-mirror-symmetry}
\end{proposition}
\begin{proof}If $\alpha$ is not very stable then there exists  $\beta\neq\alpha \in W^+_\alpha\cap \calC$  in the upward flow of $
	\alpha$. Then $\lim_{\lambda\to \infty}(\lambda\cdot \beta)\in M^\T$ is fixed by the $\T$-action. Moreover a linearized $\T$-action on a very ample line bundle on $M$ (which exists, because $M$ is normal and quasi-projective and \cite[Corollary 7.2]{dolgachev}) embeds it equivariantly in some projective space $M\subset {\mathbb P}^N$ with linear $\T$-action, where the closure of every non-trivial $\T$-orbit is linear, thus has two distinct $\T$-fixed points. Consequently, $$\lim_{\lambda\to \infty}(\lambda\cdot \beta)\neq\lim_{\lambda\to 0}(\lambda\cdot \beta)=\alpha$$and thus $W^+_\alpha$ is not closed.
	
	Assume now that $\alpha\in M^{s\T}$ is very stable i.e. $$W^+_{\alpha}\cap \calC=\{\alpha\}$$ as sets.  Then we will determine the image of  $W^+_\alpha\setminus \calC$ in the geometric quotient $$Z=(M\setminus \calC)/\T.$$ By construction it is $$\P(W^+_\alpha)=W^+_\alpha\setminus \{\alpha\}/\T\subset Z$$  isomorphic to the  weighted projective space $\P(T_\alpha^+)=(T_\alpha^+\setminus \{0\})/\T$. By the construction of
	\cite[(1.2.3)]{hausel-large} (c.f. also \cite[\S 11]{simpson3} and \cite{hausel-compact} in the Higgs moduli space case) we have the compactification $$\overline{M}=(M\times\C)/\! /\T=\left(M\times\C\setminus \calC\times\{0\}\right)/\T\cong M\sqcup Z.$$ We 
	see that the geometric quotient $$\P(W^+_\alpha\times \C)= \left((W^+_\alpha\times \C)\setminus (\alpha,0)\right)/\T= W^+_\alpha\sqcup \P(W^+_\alpha) \subset M\sqcup Z = \overline{M}$$ is also a weighted projective space, thus projective, thus closed in $\overline{M}$. It follows that all boundary points of the closure $\overline{U}_\alpha$ in $\overline{M}$ lie on the divisor at infinity $Z\subset \overline{M}\setminus M$, thus $W^+_\alpha$ is closed in $M$. 
\end{proof}
